\chapter{k-Nearest Neighbors}
\label{chap:knn}

\begin{myremark}[Algorithmusart in diesem Kapitel]
    Funktionsart: \textbf{Klassifikation} (und Regression).
    Umsetzungsart: \textbf{selbstlernend}.
\end{myremark}

\section{Grundidee: Klassifizieren durch Ähnlichkeit}

\begin{mydefinition}[k-Nearest Neighbors (KNN)]
    Beim \textbf{KNN}-Verfahren wird ein neuer Datenpunkt nicht durch ein gelerntes Modell,
    sondern durch die \textbf{$k$ ähnlichsten Punkte} (Nachbarn) im Trainingsset klassifiziert.
    Die häufigste Klasse in der Nachbarschaft bestimmt die Vorhersage.
\end{mydefinition}

\begin{myexample}
    $k = 3$: Der neue Punkt hat zwei Nachbarn der Klasse \enquote{A} und einen der Klasse \enquote{B}.
    Vorhersage: \enquote{A}.
\end{myexample}

\begin{figure}[H]
    \centering
    \foreach \k in {1, 5, 15, 50} {
            \begin{subfigure}[H]{0.45\textwidth}
                \includegraphics[width=\textwidth]{Grundlagen_Info/10_KI_und_Algorithmen/Figures/knn_boundary_\k.pdf}
                \caption{$k$ = \k}
            \end{subfigure}
        }
    \caption{KNN-Entscheidungsgrenze für verschiedene Werte von $k$}
    \label{fig:knn-boundary-15}
\end{figure}

\begin{myexercise}[KNN-Entscheidungsgrenze verstehen]
    Betrachten Sie \cref{fig:knn-boundary-15}:
    \begin{itemize}
        \item Wie verändert sich die Entscheidungsgrenze mit steigendem $k$? Wo erkennen Sie Overfitting? Wo Underfitting?
        \item Warum könnte $k=1$ zu Overfitting führen?
        \item Warum könnte ein zu grosses $k$ zu Underfitting führen?
    \end{itemize}
    \begin{myanswer}
        \begin{itemize}
            \item Mit steigendem $k$ wird die Entscheidungsgrenze glatter. $k=1$ zeigt eine sehr unregelmässige Grenze, die stark auf Ausreisser reagiert (Overfitting). $k=50$ zeigt eine sehr glatte Grenze, die möglicherweise wichtige Strukturen übersieht (Underfitting).
            \item $k=1$ führt zu Overfitting, weil jeder Punkt nur seinen nächsten Nachbarn betrachtet. Ein Ausreisser könnte also eine falsche Klasse vorhersagen.
            \item Ein zu grosses $k$ führt zu Underfitting, weil es zu viele Nachbarn berücksichtigt und dadurch die Vorhersage zu allgemein wird. Es könnte wichtige lokale Muster übersehen.
        \end{itemize}
    \end{myanswer}
\end{myexercise}

\section{Distanzmasse}

Die Ähnlichkeit wird über eine \textbf{Distanz} gemessen. Gängig ist die \textbf{euklidische Distanz}:
\[
    \begin{aligned}
        d(a, b) &= \sqrt{(x_1^a - x_1^b)^2 + (x_2^a - x_2^b)^2 + \cdots + (x_n^a - x_n^b)^2 } \\
                & = \sqrt{\sum_{i=1}^n (x_i^a - x_i^b)^2}
    \end{aligned}
\]

\begin{mycaution}
    Wenn Features sehr unterschiedliche Wertebereiche haben (z.\,B. Alter vs.\ Einkommen),
    dominiert das Feature mit grösseren Werten die Distanz.
    \textbf{Normalisierung} (vgl.  \cref{chap:datenvorbereitung}) ist daher bei KNN besonders wichtig.
\end{mycaution}

\section{Wahl von \texorpdfstring{$k$}{k}}

\begin{myexercise}[Einfluss von $k$]
    Beschreiben Sie den Unterschied zwischen kleinem und grossem $k$:
    \begin{itemize}
        \item Wann reagiert KNN sehr empfindlich auf Ausreisser?
        \item Wann wird die Entscheidungsgrenze sehr \enquote{glatt}?
        \item Was passiert bei $k = 1$? Und bei $k = N$ (alle Trainingspunkte)?
    \end{itemize}
\end{myexercise}

\begin{mychallenge}[Optimales $k$ finden]
    Testen Sie im Notebook $k = 1, 3, 5, 9, 15$ und dokumentieren Sie:
    \begin{enumerate}
        \item Trainingsgenauigkeit
        \item Testgenauigkeit
        \item Visualisierung der Entscheidungsgrenze
    \end{enumerate}
    Begründen Sie Ihre Wahl des besten $k$.
\end{mychallenge}

\section{KNN für Regression}

KNN kann auch für \textbf{Regression} eingesetzt werden:
Anstelle der häufigsten Klasse wird der \textbf{Durchschnitt} der $k$ Nachbar-Ausgabewerte
als Vorhersage verwendet.
\section{Gesellschaftliche Aspekte: Normalisierungsbias und Datenqualität}

\begin{mydefinition}[Stille Voraussetzungen von KNN]
    KNN ist einfach und intuitiv, aber macht implizite Annahmen über Datenqualität und Merkmalsdefinition. Eine Normalisierung ist nicht optional -- sie ist unvermeidbar. Die Frage ist: Wer entscheidet, wie normalisiert wird?
\end{mydefinition}

\begin{myexample}[Normalisierung als Wertsetzung]
    Betrachten Sie ein KNN-System zur Risikoeinschätzung bei Versicherungen. Merkmale: \enquote{Alter} und \enquote{Jahreseinkommen}. Min-Max-Normalisierung auf \([0,1]\):
    \[
        \text{Alter normalisiert: } \frac{30-18}{100-18} = 0.146
    \]
    \[
        \text{Einkommen normalisiert: } \frac{50000-10000}{200000-10000} = 0.235
    \]
    Nach Normalisierung haben beide Merkmale gleich \enquote{Gewicht}. Aber ist ein Jahrzehnt Alter wirklich so riskant wie 40'000 CHF Einkommensdifferenz? Diese Entscheidung ist nicht objektiv, sondern Ergebnis von Wertsetzungen.

    \textbf{Real-World-Risiko:} Wenn Einkommensunterschiede historisch mit Ethnizität oder Geschlecht korrelieren, verstärkt eine naive Normalisierung diese Disparität.
\end{myexample}

\begin{mychallenge}[Normalisierung und Fairness]
    Sie trainieren ein KNN-System zur Früherkennung von Gesundheitsrisiken in einer Klinik. Merkmale sind u.a. BMI, Blutdruck und \enquote{Anzahl Arztbesuche im letzten Jahr}.
    \begin{enumerate}
        \item Warum könnten \enquote{Arztbesuche} ein biased Feature sein?
        \item Wie könnte Normalisierung diese Verzerrung verstärken?
        \item Welche Alternative würden Sie vorschlagen -- einfach das Feature weglassen oder die Datensammlung ändern?
    \end{enumerate}
    \begin{myanswer}
        \begin{enumerate}
            \item Menschen mit höherem Einkommen haben oft besseren Zugang zu Ärzt*innen. Häufige Besuche könnten also Wohlstand abbilden, nicht Krankenheit.
            \item Wenn dieses Feature normalisiert und gleich gewichtet wird wie medizinische Parameter, wird Wohlstand direkt in eine \enquote{Risikoeinschätzung} übersetzt.
            \item Idealerweise: Datensammlung diversifizieren und Zugangsbarrieren untersuchen. Aber realistisch erfordert das Zeit. Schnelle Lösung: Feature-Wichtigkeitstests durchführen, um zu sehen, ob \enquote{Arztbesuche} auch nach Kontrolle für Wohlstand noch prädiktiv ist.
        \end{enumerate}
    \end{myanswer}
\end{mychallenge}

\begin{myexercise}[Praktische Umsetzung in Python]
    Implementieren Sie eine KNN-Klassifikation mit dem Datensatz \enquote{\texttt{learning\_style\_classifier.csv}} in \href{\notebookbaseurl04_KNN_Classification.ipynb}{\texttt{04\_KNN.ipynb}}. Erstellen Sie danach einen weiteren Versuch mit dem Datensatz \enquote{\texttt{music\_talent\_classifier.csv}} und vergleichen Sie die Ergebnisse für unterschiedliche Werte von $k$.

    \begin{myanswer}
        \href{\notebookbaseurl04_KNN_Classification_sol.ipynb}{\texttt{04\_KNN\_sol.ipynb}}
    \end{myanswer}
\end{myexercise}

